% =========================================================
% ppt/src/minggu-05.tex
% Slide pertemuan 5
% Bab VII — Format Instruksi Mesin
%
% Sub-CPMK 2.1 | CPMK02 | CPL02
% =========================================================

\documentclass[aspectratio=169,11pt]{beamer}
% =========================================================
% ppt/src/preamble-ppt.tex
% Tema, warna, paket, dan makro bersama seluruh deck slide.
%
% Dipakai oleh setiap berkas minggu-NN.tex dengan pola:
%   \documentclass[aspectratio=169,11pt]{beamer}
%   % =========================================================
% ppt/src/preamble-ppt.tex
% Tema, warna, paket, dan makro bersama seluruh deck slide.
%
% Dipakai oleh setiap berkas minggu-NN.tex dengan pola:
%   \documentclass[aspectratio=169,11pt]{beamer}
%   % =========================================================
% ppt/src/preamble-ppt.tex
% Tema, warna, paket, dan makro bersama seluruh deck slide.
%
% Dipakai oleh setiap berkas minggu-NN.tex dengan pola:
%   \documentclass[aspectratio=169,11pt]{beamer}
%   \input{preamble-ppt}
%
% Palet warna disalin dari buku_ajar/config/styles.tex agar
% identitas visual slide sama dengan buku ajar.
% =========================================================

% ---------------------------------------------------------
% 1. Paket
% ---------------------------------------------------------
\usepackage[T1]{fontenc}
\usepackage[utf8]{inputenc}
\usepackage[indonesian]{babel}
\usepackage{lmodern}
\usepackage{tikz}
\usetikzlibrary{arrows.meta, positioning, calc, shapes.geometric}
\usepackage{booktabs}
\usepackage{array}
\usepackage{amsmath}
\usepackage{amsfonts}
\usepackage{pifont}          % \ding untuk tanda checklist
\usepackage{listings}

% ---------------------------------------------------------
% 2. Palet warna OBE (identik dengan buku)
% ---------------------------------------------------------
\definecolor{obeBlue}{HTML}{1F3864}
\definecolor{obeBlueLight}{HTML}{DCE4F5}
\definecolor{obeTeal}{HTML}{0F6E6E}
\definecolor{obeTealLight}{HTML}{D9F0EE}
\definecolor{obeGreen}{HTML}{2E6B34}
\definecolor{obeGreenLight}{HTML}{DEEFDF}
\definecolor{obeAmber}{HTML}{8A5A00}
\definecolor{obeAmberLight}{HTML}{FDF1D6}
\definecolor{obeMaroon}{HTML}{7A2233}
\definecolor{obeMaroonLight}{HTML}{F7E1E5}
\definecolor{obeGray}{HTML}{4A4A4A}
\definecolor{obeGrayLight}{HTML}{F0F0F0}
\definecolor{codeBack}{HTML}{F7F7F4}
\definecolor{codeFrame}{HTML}{C9C9C0}

% ---------------------------------------------------------
% 3. Identitas mata kuliah (dipakai seluruh deck)
% ---------------------------------------------------------
\newcommand{\NamaMK}{Bahasa Pemrograman Tingkat Rendah}
\newcommand{\NamaMKShort}{Bhs.\ Pemrograman Tingkat Rendah}
\newcommand{\KodeMK}{MWCA010}
\newcommand{\Prodi}{Program Studi Teknik Informatika}
\newcommand{\Fakultas}{Fakultas Teknologi Informasi}
\newcommand{\PerguruanTinggi}{Universitas Bale Bandung}
\newcommand{\DosenPengampu}{Cecep Suwanda, S.Si., M.Kom.}
\newcommand{\SemesterTeks}{Semester Ganjil}

% \PertemuanTeks dan \JudulBab ditetapkan di masing-masing
% berkas minggu-NN.tex sebelum \begin{document}.
\providecommand{\PertemuanTeks}{Pertemuan}

% ---------------------------------------------------------
% 4. Tema beamer
% ---------------------------------------------------------
\usetheme{default}
\setbeamertemplate{navigation symbols}{}

\setbeamercolor{normal text}{fg=black!88}
\setbeamercolor{structure}{fg=obeBlue}
\setbeamercolor{alerted text}{fg=obeMaroon}

% --- judul slide: pita biru penuh lebar ---
\setbeamercolor{frametitle}{fg=white,bg=obeBlue}
\setbeamerfont{frametitle}{series=\bfseries,size=\large}
\setbeamertemplate{frametitle}{%
    \nointerlineskip%
    \begin{beamercolorbox}[%
        wd=\paperwidth,ht=2.9ex,dp=1.5ex,
        leftskip=0.4cm,rightskip=0.4cm]{frametitle}%
        \usebeamerfont{frametitle}\insertframetitle%
    \end{beamercolorbox}%
    \nointerlineskip%
    \hbox{\textcolor{obeTeal}{\rule{\paperwidth}{1.1pt}}}%
}

% --- kaki slide: identitas mata kuliah, pertemuan, nomor ---
\setbeamercolor{footline}{fg=white,bg=obeGray}
\setbeamerfont{footline}{size=\scriptsize,family=\sffamily}
\setbeamertemplate{footline}{%
    \leavevmode%
    \hbox{%
        \begin{beamercolorbox}[wd=0.42\paperwidth,ht=2.1ex,dp=1.1ex,
            leftskip=0.4cm]{footline}%
            \usebeamerfont{footline}%
            \KodeMK\ \textbullet\ \NamaMKShort%
        \end{beamercolorbox}%
        \begin{beamercolorbox}[wd=0.32\paperwidth,ht=2.1ex,dp=1.1ex,
            center]{footline}%
            \usebeamerfont{footline}\PertemuanTeks%
        \end{beamercolorbox}%
        \begin{beamercolorbox}[wd=0.26\paperwidth,ht=2.1ex,dp=1.1ex,
            rightskip=0.4cm]{footline}%
            \usebeamerfont{footline}%
            \insertframenumber\,/\,\inserttotalframenumber%
        \end{beamercolorbox}%
    }%
    \vskip0pt%
}

% --- daftar ---
\setbeamercolor{itemize item}{fg=obeTeal}
\setbeamercolor{itemize subitem}{fg=obeTeal}
\setbeamercolor{enumerate item}{fg=obeBlue}
\setbeamertemplate{itemize item}{\raise0.2ex\hbox{\tiny$\blacksquare$}}
\setbeamertemplate{itemize subitem}{\raise0.2ex\hbox{\scriptsize$\bullet$}}

% --- kotak (pengganti tcolorbox pada buku) ---
\setbeamercolor{block title}{bg=obeBlueLight,fg=obeBlue}
\setbeamercolor{block body}{bg=obeBlueLight!45,fg=black!88}
\setbeamercolor{block title alerted}{bg=obeMaroonLight,fg=obeMaroon}
\setbeamercolor{block body alerted}{bg=obeMaroonLight!45,fg=black!88}
\setbeamercolor{block title example}{bg=obeTealLight,fg=obeTeal}
\setbeamercolor{block body example}{bg=obeTealLight!45,fg=black!88}

% ---------------------------------------------------------
% 5. Makro notasi (sama persis dengan buku)
%
% Catatan: \note TIDAK didefinisikan di sini karena beamer
% sudah memakai \note untuk catatan pembicara. Untuk catatan
% di slide dipakai \catatan.
% ---------------------------------------------------------
\newcommand{\code}[1]{\texttt{#1}}
\newcommand{\asm}[1]{\texttt{#1}}
\newcommand{\reg}[1]{\texttt{#1}}
\newcommand{\flag}[1]{\texttt{#1}}
\newcommand{\dir}[1]{\texttt{#1}}
\newcommand{\file}[1]{\texttt{#1}}
\newcommand{\eng}[1]{\textit{#1}}

\newcommand{\hex}[1]{\texttt{#1h}}
\newcommand{\bin}[1]{\texttt{#1b}}
\newcommand{\dec}[1]{\texttt{#1}}

\newcommand{\important}[1]{\textbf{Penting:} #1}
\newcommand{\catatan}[1]{\textbf{Catatan:} #1}
\newcommand{\term}[1]{\textbf{#1}}

\newcommand{\CPL}[1]{\textbf{CPL#1}}
\newcommand{\CPMK}[1]{\textbf{CPMK#1}}
\newcommand{\SubCPMK}[1]{\textbf{Sub-CPMK #1}}

\newcommand{\checkbox}{\ding{113}\hspace{0.35em}}
\newcommand{\checked}{\ding{51}\hspace{0.35em}}
\newcommand{\crossed}{\ding{55}\hspace{0.35em}}

\newcommand{\flowarrow}{\ensuremath{\rightarrow}}

% ---------------------------------------------------------
% 6. Makro tata letak slide
% ---------------------------------------------------------

% Keterangan sumber materi (ketertelusuran ke buku)
\newcommand{\Sumber}[1]{%
    \vfill
    {\scriptsize\color{obeGray}Sumber: #1\par}%
}

% Frame judul deck: \JudulDeck{<judul bab>}{<pertemuan>}
\newcommand{\JudulDeck}[2]{%
\begin{frame}[plain]
    \vspace{0.25cm}
    {\color{obeBlue}\sffamily\bfseries\small \NamaMK\ (\KodeMK)\par}
    \vspace{0.12cm}
    {\color{obeGray}\sffamily\scriptsize
     \Prodi\ \textbullet\ \Fakultas\ \textbullet\ \PerguruanTinggi\par}
    \vfill
    {\color{obeBlue}\sffamily\LARGE\bfseries #1\par}
    \vspace{0.4cm}
    {\color{obeTeal}\sffamily\large #2\par}
    \vspace{0.45cm}
    {\color{obeTeal}\rule{\textwidth}{1pt}}
    \vfill
    {\color{obeGray}\sffamily\small Dosen Pengampu\par}
    \vspace{0.08cm}
    {\sffamily\bfseries \DosenPengampu\par}
    \vspace{0.2cm}
    {\color{obeGray}\sffamily\small \SemesterTeks\par}
    \vspace{0.25cm}
\end{frame}
}

% Frame pemisah bagian
\newcommand{\PemisahBagian}[1]{%
\begin{frame}[plain]
    \vfill
    \begin{center}
        {\color{obeTeal}\sffamily\small\bfseries BAGIAN\par}
        \vspace{0.35cm}
        {\color{obeBlue}\sffamily\Large\bfseries #1\par}
        \vspace{0.35cm}
        {\color{obeTeal}\rule{0.55\textwidth}{0.9pt}}
    \end{center}
    \vfill
\end{frame}
}

% ---------------------------------------------------------
% 7. Gaya penampil kode
% ---------------------------------------------------------
\input{listing-ppt}

%
% Palet warna disalin dari buku_ajar/config/styles.tex agar
% identitas visual slide sama dengan buku ajar.
% =========================================================

% ---------------------------------------------------------
% 1. Paket
% ---------------------------------------------------------
\usepackage[T1]{fontenc}
\usepackage[utf8]{inputenc}
\usepackage[indonesian]{babel}
\usepackage{lmodern}
\usepackage{tikz}
\usetikzlibrary{arrows.meta, positioning, calc, shapes.geometric}
\usepackage{booktabs}
\usepackage{array}
\usepackage{amsmath}
\usepackage{amsfonts}
\usepackage{pifont}          % \ding untuk tanda checklist
\usepackage{listings}

% ---------------------------------------------------------
% 2. Palet warna OBE (identik dengan buku)
% ---------------------------------------------------------
\definecolor{obeBlue}{HTML}{1F3864}
\definecolor{obeBlueLight}{HTML}{DCE4F5}
\definecolor{obeTeal}{HTML}{0F6E6E}
\definecolor{obeTealLight}{HTML}{D9F0EE}
\definecolor{obeGreen}{HTML}{2E6B34}
\definecolor{obeGreenLight}{HTML}{DEEFDF}
\definecolor{obeAmber}{HTML}{8A5A00}
\definecolor{obeAmberLight}{HTML}{FDF1D6}
\definecolor{obeMaroon}{HTML}{7A2233}
\definecolor{obeMaroonLight}{HTML}{F7E1E5}
\definecolor{obeGray}{HTML}{4A4A4A}
\definecolor{obeGrayLight}{HTML}{F0F0F0}
\definecolor{codeBack}{HTML}{F7F7F4}
\definecolor{codeFrame}{HTML}{C9C9C0}

% ---------------------------------------------------------
% 3. Identitas mata kuliah (dipakai seluruh deck)
% ---------------------------------------------------------
\newcommand{\NamaMK}{Bahasa Pemrograman Tingkat Rendah}
\newcommand{\NamaMKShort}{Bhs.\ Pemrograman Tingkat Rendah}
\newcommand{\KodeMK}{MWCA010}
\newcommand{\Prodi}{Program Studi Teknik Informatika}
\newcommand{\Fakultas}{Fakultas Teknologi Informasi}
\newcommand{\PerguruanTinggi}{Universitas Bale Bandung}
\newcommand{\DosenPengampu}{Cecep Suwanda, S.Si., M.Kom.}
\newcommand{\SemesterTeks}{Semester Ganjil}

% \PertemuanTeks dan \JudulBab ditetapkan di masing-masing
% berkas minggu-NN.tex sebelum \begin{document}.
\providecommand{\PertemuanTeks}{Pertemuan}

% ---------------------------------------------------------
% 4. Tema beamer
% ---------------------------------------------------------
\usetheme{default}
\setbeamertemplate{navigation symbols}{}

\setbeamercolor{normal text}{fg=black!88}
\setbeamercolor{structure}{fg=obeBlue}
\setbeamercolor{alerted text}{fg=obeMaroon}

% --- judul slide: pita biru penuh lebar ---
\setbeamercolor{frametitle}{fg=white,bg=obeBlue}
\setbeamerfont{frametitle}{series=\bfseries,size=\large}
\setbeamertemplate{frametitle}{%
    \nointerlineskip%
    \begin{beamercolorbox}[%
        wd=\paperwidth,ht=2.9ex,dp=1.5ex,
        leftskip=0.4cm,rightskip=0.4cm]{frametitle}%
        \usebeamerfont{frametitle}\insertframetitle%
    \end{beamercolorbox}%
    \nointerlineskip%
    \hbox{\textcolor{obeTeal}{\rule{\paperwidth}{1.1pt}}}%
}

% --- kaki slide: identitas mata kuliah, pertemuan, nomor ---
\setbeamercolor{footline}{fg=white,bg=obeGray}
\setbeamerfont{footline}{size=\scriptsize,family=\sffamily}
\setbeamertemplate{footline}{%
    \leavevmode%
    \hbox{%
        \begin{beamercolorbox}[wd=0.42\paperwidth,ht=2.1ex,dp=1.1ex,
            leftskip=0.4cm]{footline}%
            \usebeamerfont{footline}%
            \KodeMK\ \textbullet\ \NamaMKShort%
        \end{beamercolorbox}%
        \begin{beamercolorbox}[wd=0.32\paperwidth,ht=2.1ex,dp=1.1ex,
            center]{footline}%
            \usebeamerfont{footline}\PertemuanTeks%
        \end{beamercolorbox}%
        \begin{beamercolorbox}[wd=0.26\paperwidth,ht=2.1ex,dp=1.1ex,
            rightskip=0.4cm]{footline}%
            \usebeamerfont{footline}%
            \insertframenumber\,/\,\inserttotalframenumber%
        \end{beamercolorbox}%
    }%
    \vskip0pt%
}

% --- daftar ---
\setbeamercolor{itemize item}{fg=obeTeal}
\setbeamercolor{itemize subitem}{fg=obeTeal}
\setbeamercolor{enumerate item}{fg=obeBlue}
\setbeamertemplate{itemize item}{\raise0.2ex\hbox{\tiny$\blacksquare$}}
\setbeamertemplate{itemize subitem}{\raise0.2ex\hbox{\scriptsize$\bullet$}}

% --- kotak (pengganti tcolorbox pada buku) ---
\setbeamercolor{block title}{bg=obeBlueLight,fg=obeBlue}
\setbeamercolor{block body}{bg=obeBlueLight!45,fg=black!88}
\setbeamercolor{block title alerted}{bg=obeMaroonLight,fg=obeMaroon}
\setbeamercolor{block body alerted}{bg=obeMaroonLight!45,fg=black!88}
\setbeamercolor{block title example}{bg=obeTealLight,fg=obeTeal}
\setbeamercolor{block body example}{bg=obeTealLight!45,fg=black!88}

% ---------------------------------------------------------
% 5. Makro notasi (sama persis dengan buku)
%
% Catatan: \note TIDAK didefinisikan di sini karena beamer
% sudah memakai \note untuk catatan pembicara. Untuk catatan
% di slide dipakai \catatan.
% ---------------------------------------------------------
\newcommand{\code}[1]{\texttt{#1}}
\newcommand{\asm}[1]{\texttt{#1}}
\newcommand{\reg}[1]{\texttt{#1}}
\newcommand{\flag}[1]{\texttt{#1}}
\newcommand{\dir}[1]{\texttt{#1}}
\newcommand{\file}[1]{\texttt{#1}}
\newcommand{\eng}[1]{\textit{#1}}

\newcommand{\hex}[1]{\texttt{#1h}}
\newcommand{\bin}[1]{\texttt{#1b}}
\newcommand{\dec}[1]{\texttt{#1}}

\newcommand{\important}[1]{\textbf{Penting:} #1}
\newcommand{\catatan}[1]{\textbf{Catatan:} #1}
\newcommand{\term}[1]{\textbf{#1}}

\newcommand{\CPL}[1]{\textbf{CPL#1}}
\newcommand{\CPMK}[1]{\textbf{CPMK#1}}
\newcommand{\SubCPMK}[1]{\textbf{Sub-CPMK #1}}

\newcommand{\checkbox}{\ding{113}\hspace{0.35em}}
\newcommand{\checked}{\ding{51}\hspace{0.35em}}
\newcommand{\crossed}{\ding{55}\hspace{0.35em}}

\newcommand{\flowarrow}{\ensuremath{\rightarrow}}

% ---------------------------------------------------------
% 6. Makro tata letak slide
% ---------------------------------------------------------

% Keterangan sumber materi (ketertelusuran ke buku)
\newcommand{\Sumber}[1]{%
    \vfill
    {\scriptsize\color{obeGray}Sumber: #1\par}%
}

% Frame judul deck: \JudulDeck{<judul bab>}{<pertemuan>}
\newcommand{\JudulDeck}[2]{%
\begin{frame}[plain]
    \vspace{0.25cm}
    {\color{obeBlue}\sffamily\bfseries\small \NamaMK\ (\KodeMK)\par}
    \vspace{0.12cm}
    {\color{obeGray}\sffamily\scriptsize
     \Prodi\ \textbullet\ \Fakultas\ \textbullet\ \PerguruanTinggi\par}
    \vfill
    {\color{obeBlue}\sffamily\LARGE\bfseries #1\par}
    \vspace{0.4cm}
    {\color{obeTeal}\sffamily\large #2\par}
    \vspace{0.45cm}
    {\color{obeTeal}\rule{\textwidth}{1pt}}
    \vfill
    {\color{obeGray}\sffamily\small Dosen Pengampu\par}
    \vspace{0.08cm}
    {\sffamily\bfseries \DosenPengampu\par}
    \vspace{0.2cm}
    {\color{obeGray}\sffamily\small \SemesterTeks\par}
    \vspace{0.25cm}
\end{frame}
}

% Frame pemisah bagian
\newcommand{\PemisahBagian}[1]{%
\begin{frame}[plain]
    \vfill
    \begin{center}
        {\color{obeTeal}\sffamily\small\bfseries BAGIAN\par}
        \vspace{0.35cm}
        {\color{obeBlue}\sffamily\Large\bfseries #1\par}
        \vspace{0.35cm}
        {\color{obeTeal}\rule{0.55\textwidth}{0.9pt}}
    \end{center}
    \vfill
\end{frame}
}

% ---------------------------------------------------------
% 7. Gaya penampil kode
% ---------------------------------------------------------
% =========================================================
% ppt/src/listing-ppt.tex
% Gaya penampil kode untuk slide.
%
% Diturunkan dari buku_ajar/config/listings.tex, tetapi
% disesuaikan untuk slide: ukuran lebih kecil (karena frame
% 16:9 dan kode harus utuh terlihat) dan jarak lebih rapat.
% =========================================================

% ---------------------------------------------------------
% 1. Definisi bahasa MASM/TASM 8086 (salinan dari buku)
% ---------------------------------------------------------
\lstdefinelanguage{MASM}{
    morekeywords={
        MOV, XCHG, LEA, LDS, LES, PUSH, POP, PUSHA, POPA, PUSHF, POPF,
        LAHF, SAHF, XLAT, IN, OUT, CBW, CWD,
        ADD, ADC, SUB, SBB, INC, DEC, NEG, CMP, MUL, IMUL, DIV, IDIV,
        AAA, AAS, AAM, AAD, DAA, DAS,
        AND, OR, XOR, NOT, TEST, SHL, SHR, SAL, SAR, ROL, ROR, RCL, RCR,
        JMP, JE, JZ, JNE, JNZ, JG, JGE, JL, JLE, JA, JAE, JB, JBE,
        JC, JNC, JS, JNS, JO, JNO, JP, JPE, JNP, JPO, JCXZ,
        LOOP, LOOPE, LOOPZ, LOOPNE, LOOPNZ, CALL, RET, INT, IRET, INTO,
        HLT, NOP, CLI, STI, CLD, STD, CLC, STC, CMC, WAIT, ESC, LOCK,
        DB, DW, DD, DQ, DT, EQU, ORG, SEGMENT, ENDS, ASSUME, END,
        PROC, ENDP, MACRO, ENDM, LOCAL, PUBLIC, EXTRN, EXTERN,
        INCLUDE, INCLUDELIB, MODEL, STACK, CODE, DATA, CONST, ALIGN, EVEN,
        OFFSET, PTR, BYTE, WORD, DWORD, QWORD, TBYTE,
        NEAR, FAR, SHORT, LABEL, TYPE, SIZE, LENGTH, DUP, THIS,
        AX, BX, CX, DX, AL, AH, BL, BH, CL, CH, DL, DH,
        SP, BP, SI, DI, IP, CS, DS, SS, ES, FLAGS,
        CF, PF, AF, ZF, TF, IF, DF, OF
    },
    sensitive=false,
    morecomment=[l]{;},
    morestring=[b]',
    morestring=[b]"
}

% ---------------------------------------------------------
% 2. Gaya kode assembly untuk slide
% ---------------------------------------------------------
\lstdefinestyle{asmslide}{
    language=MASM,
    basicstyle=\ttfamily\scriptsize,
    keywordstyle=\color{obeBlue}\bfseries,
    commentstyle=\color{obeGreen}\itshape,
    stringstyle=\color{obeMaroon},
    numberstyle=\tiny\color{obeGray},
    numbers=left,
    numbersep=5pt,
    stepnumber=1,
    backgroundcolor=\color{codeBack},
    frame=single,
    rulecolor=\color{codeFrame},
    framesep=4pt,
    breaklines=true,
    breakatwhitespace=false,
    % CATATAN PENTING: opsi `postbreak` TIDAK boleh dipakai di sini.
    % Di beamer, `breaklines=true` bersama `postbreak=<kode>` membuat
    % listings gagal dirakit dengan galat
    %   "Missing number, treated as zero <to be read again> \kern".
    % Karena \lstset{style=asmslide} bersifat global, opsi postbreak
    % juga bocor ke gaya lain (mis. outslide/keluaran) dan merusaknya.
    % Di buku ajar (kelas `book`) justru sebaliknya: postbreak aman dan
    % dipakai. Jadi larangan ini khusus untuk slide.
    showstringspaces=false,
    tabsize=4,
    columns=fullflexible,
    keepspaces=true,
    upquote=true,
    xleftmargin=12pt,
    aboveskip=0.4em,
    belowskip=0.4em
}

% ---------------------------------------------------------
% 3. Gaya keluaran program / sesi debugger untuk slide
% ---------------------------------------------------------
\lstdefinestyle{outslide}{
    basicstyle=\ttfamily\scriptsize,
    keywordstyle=\color{black},
    commentstyle=\color{obeGray},
    stringstyle=\color{black},
    numbers=none,
    backgroundcolor=\color{obeGrayLight},
    frame=single,
    rulecolor=\color{codeFrame},
    framesep=4pt,
    breaklines=true,
    showstringspaces=false,
    columns=fullflexible,
    keepspaces=true,
    upquote=true,
    xleftmargin=4pt,
    aboveskip=0.4em,
    belowskip=0.4em
}

\lstset{style=asmslide}

% ---------------------------------------------------------
% 4. Environment ringkas untuk dipakai di dalam frame
%
% Frame yang memuat kedua environment ini WAJIB memakai
% opsi [fragile], karena isinya verbatim.
% ---------------------------------------------------------
\lstnewenvironment{kodeasm}[1][]{%
    \lstset{style=asmslide, #1}%
}{}

\lstnewenvironment{keluaran}[1][]{%
    \lstset{style=outslide, #1}%
}{}


%
% Palet warna disalin dari buku_ajar/config/styles.tex agar
% identitas visual slide sama dengan buku ajar.
% =========================================================

% ---------------------------------------------------------
% 1. Paket
% ---------------------------------------------------------
\usepackage[T1]{fontenc}
\usepackage[utf8]{inputenc}
\usepackage[indonesian]{babel}
\usepackage{lmodern}
\usepackage{tikz}
\usetikzlibrary{arrows.meta, positioning, calc, shapes.geometric}
\usepackage{booktabs}
\usepackage{array}
\usepackage{amsmath}
\usepackage{amsfonts}
\usepackage{pifont}          % \ding untuk tanda checklist
\usepackage{listings}

% ---------------------------------------------------------
% 2. Palet warna OBE (identik dengan buku)
% ---------------------------------------------------------
\definecolor{obeBlue}{HTML}{1F3864}
\definecolor{obeBlueLight}{HTML}{DCE4F5}
\definecolor{obeTeal}{HTML}{0F6E6E}
\definecolor{obeTealLight}{HTML}{D9F0EE}
\definecolor{obeGreen}{HTML}{2E6B34}
\definecolor{obeGreenLight}{HTML}{DEEFDF}
\definecolor{obeAmber}{HTML}{8A5A00}
\definecolor{obeAmberLight}{HTML}{FDF1D6}
\definecolor{obeMaroon}{HTML}{7A2233}
\definecolor{obeMaroonLight}{HTML}{F7E1E5}
\definecolor{obeGray}{HTML}{4A4A4A}
\definecolor{obeGrayLight}{HTML}{F0F0F0}
\definecolor{codeBack}{HTML}{F7F7F4}
\definecolor{codeFrame}{HTML}{C9C9C0}

% ---------------------------------------------------------
% 3. Identitas mata kuliah (dipakai seluruh deck)
% ---------------------------------------------------------
\newcommand{\NamaMK}{Bahasa Pemrograman Tingkat Rendah}
\newcommand{\NamaMKShort}{Bhs.\ Pemrograman Tingkat Rendah}
\newcommand{\KodeMK}{MWCA010}
\newcommand{\Prodi}{Program Studi Teknik Informatika}
\newcommand{\Fakultas}{Fakultas Teknologi Informasi}
\newcommand{\PerguruanTinggi}{Universitas Bale Bandung}
\newcommand{\DosenPengampu}{Cecep Suwanda, S.Si., M.Kom.}
\newcommand{\SemesterTeks}{Semester Ganjil}

% \PertemuanTeks dan \JudulBab ditetapkan di masing-masing
% berkas minggu-NN.tex sebelum \begin{document}.
\providecommand{\PertemuanTeks}{Pertemuan}

% ---------------------------------------------------------
% 4. Tema beamer
% ---------------------------------------------------------
\usetheme{default}
\setbeamertemplate{navigation symbols}{}

\setbeamercolor{normal text}{fg=black!88}
\setbeamercolor{structure}{fg=obeBlue}
\setbeamercolor{alerted text}{fg=obeMaroon}

% --- judul slide: pita biru penuh lebar ---
\setbeamercolor{frametitle}{fg=white,bg=obeBlue}
\setbeamerfont{frametitle}{series=\bfseries,size=\large}
\setbeamertemplate{frametitle}{%
    \nointerlineskip%
    \begin{beamercolorbox}[%
        wd=\paperwidth,ht=2.9ex,dp=1.5ex,
        leftskip=0.4cm,rightskip=0.4cm]{frametitle}%
        \usebeamerfont{frametitle}\insertframetitle%
    \end{beamercolorbox}%
    \nointerlineskip%
    \hbox{\textcolor{obeTeal}{\rule{\paperwidth}{1.1pt}}}%
}

% --- kaki slide: identitas mata kuliah, pertemuan, nomor ---
\setbeamercolor{footline}{fg=white,bg=obeGray}
\setbeamerfont{footline}{size=\scriptsize,family=\sffamily}
\setbeamertemplate{footline}{%
    \leavevmode%
    \hbox{%
        \begin{beamercolorbox}[wd=0.42\paperwidth,ht=2.1ex,dp=1.1ex,
            leftskip=0.4cm]{footline}%
            \usebeamerfont{footline}%
            \KodeMK\ \textbullet\ \NamaMKShort%
        \end{beamercolorbox}%
        \begin{beamercolorbox}[wd=0.32\paperwidth,ht=2.1ex,dp=1.1ex,
            center]{footline}%
            \usebeamerfont{footline}\PertemuanTeks%
        \end{beamercolorbox}%
        \begin{beamercolorbox}[wd=0.26\paperwidth,ht=2.1ex,dp=1.1ex,
            rightskip=0.4cm]{footline}%
            \usebeamerfont{footline}%
            \insertframenumber\,/\,\inserttotalframenumber%
        \end{beamercolorbox}%
    }%
    \vskip0pt%
}

% --- daftar ---
\setbeamercolor{itemize item}{fg=obeTeal}
\setbeamercolor{itemize subitem}{fg=obeTeal}
\setbeamercolor{enumerate item}{fg=obeBlue}
\setbeamertemplate{itemize item}{\raise0.2ex\hbox{\tiny$\blacksquare$}}
\setbeamertemplate{itemize subitem}{\raise0.2ex\hbox{\scriptsize$\bullet$}}

% --- kotak (pengganti tcolorbox pada buku) ---
\setbeamercolor{block title}{bg=obeBlueLight,fg=obeBlue}
\setbeamercolor{block body}{bg=obeBlueLight!45,fg=black!88}
\setbeamercolor{block title alerted}{bg=obeMaroonLight,fg=obeMaroon}
\setbeamercolor{block body alerted}{bg=obeMaroonLight!45,fg=black!88}
\setbeamercolor{block title example}{bg=obeTealLight,fg=obeTeal}
\setbeamercolor{block body example}{bg=obeTealLight!45,fg=black!88}

% ---------------------------------------------------------
% 5. Makro notasi (sama persis dengan buku)
%
% Catatan: \note TIDAK didefinisikan di sini karena beamer
% sudah memakai \note untuk catatan pembicara. Untuk catatan
% di slide dipakai \catatan.
% ---------------------------------------------------------
\newcommand{\code}[1]{\texttt{#1}}
\newcommand{\asm}[1]{\texttt{#1}}
\newcommand{\reg}[1]{\texttt{#1}}
\newcommand{\flag}[1]{\texttt{#1}}
\newcommand{\dir}[1]{\texttt{#1}}
\newcommand{\file}[1]{\texttt{#1}}
\newcommand{\eng}[1]{\textit{#1}}

\newcommand{\hex}[1]{\texttt{#1h}}
\newcommand{\bin}[1]{\texttt{#1b}}
\newcommand{\dec}[1]{\texttt{#1}}

\newcommand{\important}[1]{\textbf{Penting:} #1}
\newcommand{\catatan}[1]{\textbf{Catatan:} #1}
\newcommand{\term}[1]{\textbf{#1}}

\newcommand{\CPL}[1]{\textbf{CPL#1}}
\newcommand{\CPMK}[1]{\textbf{CPMK#1}}
\newcommand{\SubCPMK}[1]{\textbf{Sub-CPMK #1}}

\newcommand{\checkbox}{\ding{113}\hspace{0.35em}}
\newcommand{\checked}{\ding{51}\hspace{0.35em}}
\newcommand{\crossed}{\ding{55}\hspace{0.35em}}

\newcommand{\flowarrow}{\ensuremath{\rightarrow}}

% ---------------------------------------------------------
% 6. Makro tata letak slide
% ---------------------------------------------------------

% Keterangan sumber materi (ketertelusuran ke buku)
\newcommand{\Sumber}[1]{%
    \vfill
    {\scriptsize\color{obeGray}Sumber: #1\par}%
}

% Frame judul deck: \JudulDeck{<judul bab>}{<pertemuan>}
\newcommand{\JudulDeck}[2]{%
\begin{frame}[plain]
    \vspace{0.25cm}
    {\color{obeBlue}\sffamily\bfseries\small \NamaMK\ (\KodeMK)\par}
    \vspace{0.12cm}
    {\color{obeGray}\sffamily\scriptsize
     \Prodi\ \textbullet\ \Fakultas\ \textbullet\ \PerguruanTinggi\par}
    \vfill
    {\color{obeBlue}\sffamily\LARGE\bfseries #1\par}
    \vspace{0.4cm}
    {\color{obeTeal}\sffamily\large #2\par}
    \vspace{0.45cm}
    {\color{obeTeal}\rule{\textwidth}{1pt}}
    \vfill
    {\color{obeGray}\sffamily\small Dosen Pengampu\par}
    \vspace{0.08cm}
    {\sffamily\bfseries \DosenPengampu\par}
    \vspace{0.2cm}
    {\color{obeGray}\sffamily\small \SemesterTeks\par}
    \vspace{0.25cm}
\end{frame}
}

% Frame pemisah bagian
\newcommand{\PemisahBagian}[1]{%
\begin{frame}[plain]
    \vfill
    \begin{center}
        {\color{obeTeal}\sffamily\small\bfseries BAGIAN\par}
        \vspace{0.35cm}
        {\color{obeBlue}\sffamily\Large\bfseries #1\par}
        \vspace{0.35cm}
        {\color{obeTeal}\rule{0.55\textwidth}{0.9pt}}
    \end{center}
    \vfill
\end{frame}
}

% ---------------------------------------------------------
% 7. Gaya penampil kode
% ---------------------------------------------------------
% =========================================================
% ppt/src/listing-ppt.tex
% Gaya penampil kode untuk slide.
%
% Diturunkan dari buku_ajar/config/listings.tex, tetapi
% disesuaikan untuk slide: ukuran lebih kecil (karena frame
% 16:9 dan kode harus utuh terlihat) dan jarak lebih rapat.
% =========================================================

% ---------------------------------------------------------
% 1. Definisi bahasa MASM/TASM 8086 (salinan dari buku)
% ---------------------------------------------------------
\lstdefinelanguage{MASM}{
    morekeywords={
        MOV, XCHG, LEA, LDS, LES, PUSH, POP, PUSHA, POPA, PUSHF, POPF,
        LAHF, SAHF, XLAT, IN, OUT, CBW, CWD,
        ADD, ADC, SUB, SBB, INC, DEC, NEG, CMP, MUL, IMUL, DIV, IDIV,
        AAA, AAS, AAM, AAD, DAA, DAS,
        AND, OR, XOR, NOT, TEST, SHL, SHR, SAL, SAR, ROL, ROR, RCL, RCR,
        JMP, JE, JZ, JNE, JNZ, JG, JGE, JL, JLE, JA, JAE, JB, JBE,
        JC, JNC, JS, JNS, JO, JNO, JP, JPE, JNP, JPO, JCXZ,
        LOOP, LOOPE, LOOPZ, LOOPNE, LOOPNZ, CALL, RET, INT, IRET, INTO,
        HLT, NOP, CLI, STI, CLD, STD, CLC, STC, CMC, WAIT, ESC, LOCK,
        DB, DW, DD, DQ, DT, EQU, ORG, SEGMENT, ENDS, ASSUME, END,
        PROC, ENDP, MACRO, ENDM, LOCAL, PUBLIC, EXTRN, EXTERN,
        INCLUDE, INCLUDELIB, MODEL, STACK, CODE, DATA, CONST, ALIGN, EVEN,
        OFFSET, PTR, BYTE, WORD, DWORD, QWORD, TBYTE,
        NEAR, FAR, SHORT, LABEL, TYPE, SIZE, LENGTH, DUP, THIS,
        AX, BX, CX, DX, AL, AH, BL, BH, CL, CH, DL, DH,
        SP, BP, SI, DI, IP, CS, DS, SS, ES, FLAGS,
        CF, PF, AF, ZF, TF, IF, DF, OF
    },
    sensitive=false,
    morecomment=[l]{;},
    morestring=[b]',
    morestring=[b]"
}

% ---------------------------------------------------------
% 2. Gaya kode assembly untuk slide
% ---------------------------------------------------------
\lstdefinestyle{asmslide}{
    language=MASM,
    basicstyle=\ttfamily\scriptsize,
    keywordstyle=\color{obeBlue}\bfseries,
    commentstyle=\color{obeGreen}\itshape,
    stringstyle=\color{obeMaroon},
    numberstyle=\tiny\color{obeGray},
    numbers=left,
    numbersep=5pt,
    stepnumber=1,
    backgroundcolor=\color{codeBack},
    frame=single,
    rulecolor=\color{codeFrame},
    framesep=4pt,
    breaklines=true,
    breakatwhitespace=false,
    % CATATAN PENTING: opsi `postbreak` TIDAK boleh dipakai di sini.
    % Di beamer, `breaklines=true` bersama `postbreak=<kode>` membuat
    % listings gagal dirakit dengan galat
    %   "Missing number, treated as zero <to be read again> \kern".
    % Karena \lstset{style=asmslide} bersifat global, opsi postbreak
    % juga bocor ke gaya lain (mis. outslide/keluaran) dan merusaknya.
    % Di buku ajar (kelas `book`) justru sebaliknya: postbreak aman dan
    % dipakai. Jadi larangan ini khusus untuk slide.
    showstringspaces=false,
    tabsize=4,
    columns=fullflexible,
    keepspaces=true,
    upquote=true,
    xleftmargin=12pt,
    aboveskip=0.4em,
    belowskip=0.4em
}

% ---------------------------------------------------------
% 3. Gaya keluaran program / sesi debugger untuk slide
% ---------------------------------------------------------
\lstdefinestyle{outslide}{
    basicstyle=\ttfamily\scriptsize,
    keywordstyle=\color{black},
    commentstyle=\color{obeGray},
    stringstyle=\color{black},
    numbers=none,
    backgroundcolor=\color{obeGrayLight},
    frame=single,
    rulecolor=\color{codeFrame},
    framesep=4pt,
    breaklines=true,
    showstringspaces=false,
    columns=fullflexible,
    keepspaces=true,
    upquote=true,
    xleftmargin=4pt,
    aboveskip=0.4em,
    belowskip=0.4em
}

\lstset{style=asmslide}

% ---------------------------------------------------------
% 4. Environment ringkas untuk dipakai di dalam frame
%
% Frame yang memuat kedua environment ini WAJIB memakai
% opsi [fragile], karena isinya verbatim.
% ---------------------------------------------------------
\lstnewenvironment{kodeasm}[1][]{%
    \lstset{style=asmslide, #1}%
}{}

\lstnewenvironment{keluaran}[1][]{%
    \lstset{style=outslide, #1}%
}{}



\renewcommand{\PertemuanTeks}{Pertemuan 5}

\begin{document}

% =========================================================
% 1. JUDUL
% =========================================================
\JudulDeck{Format Instruksi Mesin}{Pertemuan 5}

% =========================================================
% 2. CAPAIAN DAN TUJUAN
% =========================================================
\begin{frame}{Capaian dan Tujuan Pembelajaran}
\begin{block}{Sub-CPMK 2.1}
\small Menjelaskan format instruksi mesin (\eng{opcode}, byte ModR/M,
\eng{displacement}/\eng{immediate}) dan mengaitkannya dengan operasi
yang dieksekusi perangkat keras.
\end{block}

\vspace{0.15em}
{\small\textbf{CPMK yang didukung:} \CPMK{02}\quad\textbullet\quad
\textbf{CPL yang didukung:} \CPL{02}\par}

\vspace{0.3em}
\begin{exampleblock}{Setelah pertemuan ini, mahasiswa mampu:}
\footnotesize
\begin{enumerate}
    \item menjelaskan hubungan \eng{mnemonic} assembly dan kode mesin
          yang dihasilkan assembler;
    \item menjelaskan sebab satu \eng{mnemonic} memiliki beberapa
          bentuk \eng{encoding};
    \item menguraikan susunan field instruksi 8086 beserta urutannya;
    \item memecah byte ModR/M menjadi field \code{mod},
          \code{reg/opcode}, dan \code{r/m};
    \item melakukan \eng{encoding} manual dan membacanya kembali
          dengan \eng{disassembler};
    \item menilai pengaruh format instruksi terhadap ukuran kode dan
          kecepatan eksekusi.
\end{enumerate}
\end{exampleblock}
\end{frame}

% =========================================================
% 3. PETA MATERI
% =========================================================
\begin{frame}{Peta Materi}
\small
\begin{enumerate}
    \item \textbf{Dari Mnemonic ke Kode Mesin}\\
          {\footnotesize assembly sebagai simbol kode mesin; sebab
          satu mnemonic punya banyak \eng{encoding}; format umum
          lima field; variasi panjang 1--6 byte;
          \eng{little-endian}}

    \vspace{0.25em}
    \item \textbf{Byte ModR/M, Displacement, dan Immediate}\\
          {\footnotesize field \code{mod}, \code{reg/opcode}, dan
          \code{r/m}; arti nilai \code{mod} dan \code{r/m}; dua
          pengecualian; peran \eng{displacement} dan
          \eng{immediate}; pengaruh terhadap ukuran dan kecepatan}

    \vspace{0.25em}
    \item \textbf{Encoding Manual dan Disassembler}\\
          {\footnotesize lima langkah \eng{encoding}; tabel
          \eng{opcode} umum; perintah \code{U} pada \file{DEBUG};
          keterbatasan \eng{disassembly}}
\end{enumerate}

\vspace{0.35em}
\begin{block}{Praktik pada pertemuan ini}
\small Menguraikan byte kode mesin lima instruksi secara manual, lalu
membandingkannya dengan hasil perintah \code{U} dan menghitung selisih
ukuran dua bentuk penulisan.
\end{block}
\end{frame}

% =========================================================
% BAGIAN 1
% =========================================================
\PemisahBagian{Dari Mnemonic ke Kode Mesin}

% ---------------------------------------------------------
\begin{frame}[fragile]{Assembly adalah Simbol dari Kode Mesin}
\small
\begin{itemize}
    \item Perangkat keras hanya mengenal deretan bit. Instruksi
          memindahkan isi \reg{BX} ke \reg{AX} tidak pernah tersimpan
          sebagai huruf \code{MOV}; yang ada di memori adalah dua byte
          \hex{89 D8}.

    \vspace{0.25em}
    \item Kata pendek yang mewakili operasi tersebut disebut
          \term{mnemonic}. Assembler menerjemahkan setiap
          \eng{mnemonic} beserta operandnya menjadi deretan byte kode
          mesin yang ekuivalen.

    \vspace{0.25em}
    \item \asm{MOV} tidak selalu menghasilkan \eng{opcode} yang sama:
          setiap bentuk operand memerlukan pengodean yang berbeda.
\end{itemize}

\vspace{0.35em}
\begin{keluaran}
MOV AX, BX      <- yang ditulis programmer (mnemonic)
89 D8           <- yang tersimpan di memori (kode mesin)
\end{keluaran}

\Sumber{Bab VII, Assembly adalah Simbol dari Kode Mesin}
\end{frame}

% ---------------------------------------------------------
\begin{frame}{Mengapa Satu Mnemonic Memiliki Beberapa Encoding}
\small Empat sebab utama:
\vspace{0.25em}
\begin{enumerate}
    \item \textbf{Jenis operand.} Register, sel memori, dan nilai
          segera adalah tiga jenis berbeda. \code{MOV AX, BX} memakai
          \eng{opcode} \hex{89}, sedangkan \code{MOV AX, 1234h}
          memakai \hex{B8}.

    \vspace{0.2em}
    \item \textbf{Ukuran operand.} \code{MOV AL, 5} dan
          \code{MOV AX, 5} beroperasi pada ukuran berbeda, sehingga
          \eng{opcode}-nya \hex{B0} dan \hex{B8}.

    \vspace{0.2em}
    \item \textbf{Arah perpindahan.} \hex{8B} berarti
          ``register $\leftarrow$ memori'', \hex{89} berarti
          ``memori $\leftarrow$ register''.

    \vspace{0.2em}
    \item \textbf{Mode pengalamatan.} Perbedaan alamat operand diatur
          lewat byte ModR/M dan \eng{displacement}, bukan dengan
          mengganti \eng{opcode}.
\end{enumerate}

\vspace{0.25em}
{\footnotesize\color{obeGray}Akibatnya \asm{MOV} memiliki puluhan
bentuk mesin. Programmer tidak perlu menghafalnya, tetapi pembaca hasil
\eng{disassembly} harus mengenali bentuk yang dihadapi.\par}
\end{frame}

% ---------------------------------------------------------
\begin{frame}{Format Umum Instruksi 8086}
\begin{center}
\begin{tikzpicture}[
    font=\small,
    bidang/.style={
        draw=obeBlue, line width=0.8pt, minimum height=1.45cm,
        align=center, inner sep=3pt
    }
]
    \node[bidang, fill=obeAmberLight, minimum width=1.7cm] (p)
        at (0,0) {\textbf{Prefix}\\{\scriptsize 0--1 B}\\
        {\tiny\itshape opsional}};
    \node[bidang, fill=obeBlueLight, minimum width=2.3cm,
        right=0.12cm of p] (o)
        {\textbf{Opcode}\\{\scriptsize 1--2 B}\\
        {\tiny\itshape wajib}};
    \node[bidang, fill=obeTealLight, minimum width=2.3cm,
        right=0.12cm of o] (m)
        {\textbf{ModR/M}\\{\scriptsize 0--1 B}\\
        {\tiny\itshape bila ada operand}};
    \node[bidang, fill=obeGreenLight, minimum width=2.4cm,
        right=0.12cm of m] (d)
        {\textbf{Displacement}\\{\scriptsize 0--2 B}\\
        {\tiny\itshape bila alamat bergeser}};
    \node[bidang, fill=obeMaroonLight, minimum width=2.2cm,
        right=0.12cm of d] (i)
        {\textbf{Immediate}\\{\scriptsize 0--2 B}\\
        {\tiny\itshape bila operand angka}};
\end{tikzpicture}
\end{center}

\vspace{0.4em}
{\small Urutan kelima field bersifat \textbf{tetap}; field yang tidak
diperlukan tidak ditulis sama sekali, sehingga panjang instruksi
bervariasi antara 1 sampai 6 byte.\par}

\vspace{0.3em}
{\footnotesize\color{obeGray}Contoh pemakaian penuh:
\code{MOV WORD PTR [BX+SI+1234h], 5678h} menghasilkan
\hex{C7 80 34 12 78 56} --- enam byte, semua field terpakai.\par}
\Sumber{Bab VII, Format Umum Instruksi 8086}
\end{frame}

% ---------------------------------------------------------
\begin{frame}{Ukuran Instruksi Tidak Seragam}
\footnotesize
\begin{center}
\begin{tabular}{@{}p{3.4cm}p{2.6cm}cp{4.4cm}@{}}
    \toprule
    \textbf{Instruksi} & \textbf{Kode mesin} & \textbf{Byte} &
    \textbf{Field yang dipakai} \\
    \midrule
    \asm{NOP} & \hex{90} & 1 & Opcode \\
    \addlinespace
    \asm{PUSH AX} & \hex{50} & 1 & Opcode \\
    \addlinespace
    \asm{MOV AL, 5} & \hex{B0 05} & 2 & Opcode, immediate \\
    \addlinespace
    \asm{MOV AX, BX} & \hex{89 D8} & 2 & Opcode, ModR/M \\
    \addlinespace
    \asm{MOV AX, 1234h} & \hex{B8 34 12} & 3 &
    Opcode, immediate 16 bit \\
    \addlinespace
    \asm{ADD AX, [SI+4]} & \hex{03 44 04} & 3 &
    Opcode, ModR/M, displacement \\
    \addlinespace
    \asm{MOV WORD PTR} & \hex{C7 80 34} & & Opcode, ModR/M, \\
    \asm{[BX+SI+1234h], 5678h} & \hex{12 78 56} & 6 &
    displacement, immediate \\
    \bottomrule
\end{tabular}
\end{center}

\vspace{0.25em}
{\footnotesize\color{obeGray}Intel membatasi panjang maksimum
instruksi pada 6 byte; angka inilah yang menjadi alasan antrean
pengambilan awal (\eng{prefetch queue}) 8086 berukuran 6 byte.\par}
\end{frame}

% ---------------------------------------------------------
\begin{frame}[fragile]{Membaca Kode Mesin dari Kiri ke Kanan}
\small Byte instruksi dieksekusi menurut urutan alamat, dari byte
pertama sampai terakhir. Namun nilai 16 bit di dalamnya disimpan
\term{little-endian}: byte bernilai rendah lebih dahulu.

\vspace{0.35em}
\begin{columns}[t]
\begin{column}{0.48\textwidth}
\begin{kodeasm}[numbers=none]
MOV AX, 1234h    ; nilai immediate
\end{kodeasm}
\end{column}
\begin{column}{0.46\textwidth}
{\color{obeTeal}\bfseries Byte di memori}\\[0.25em]
\begin{keluaran}
B8   34   12
\end{keluaran}
{\scriptsize byte \hex{34} lebih dahulu daripada \hex{12}}
\end{column}
\end{columns}

\vspace{0.45em}
\begin{alertblock}{Kekeliruan yang paling sering terjadi}
\small Membaca deretan \hex{34 12} sebagai \hex{3412} mengubah
nilainya, bukan lagi \hex{1234}. Kesalahan urutan ini termasuk yang
paling sering muncul pada analisis kode mesin.
\end{alertblock}
\end{frame}

% ---------------------------------------------------------
\begin{frame}{Rangkuman Bagian 1}
\small
\begin{itemize}
    \item \eng{Mnemonic} assembly hanyalah simbol bagi kode mesin.
          Yang benar-benar dibaca prosesor adalah deretan byte.

    \vspace{0.25em}
    \item Satu \eng{mnemonic} memiliki banyak bentuk \eng{encoding}
          karena jenis operand, ukuran operand, arah perpindahan data,
          dan mode pengalamatan.

    \vspace{0.25em}
    \item Format instruksi 8086 tersusun dari lima field berurutan
          tetap: awalan, \eng{opcode}, ModR/M, \eng{displacement},
          dan \eng{immediate}.

    \vspace{0.25em}
    \item Field yang tidak diperlukan tidak ditulis, sehingga panjang
          instruksi berkisar 1 sampai 6 byte.

    \vspace{0.25em}
    \item Setiap nilai 16 bit disimpan \eng{little-endian}: byte
          bernilai rendah lebih dahulu.
\end{itemize}
\end{frame}

% =========================================================
% BAGIAN 2
% =========================================================
\PemisahBagian{Byte ModR/M, Displacement, dan Immediate}

% ---------------------------------------------------------
\begin{frame}{Struktur Byte ModR/M}
\small Kehadiran byte ModR/M membuat satu \eng{opcode} dapat melayani
banyak kombinasi operand tanpa \eng{opcode} baru untuk setiap
kombinasi.

\vspace{0.4em}
\begin{center}
\begin{tikzpicture}[
    font=\small,
    bit/.style={
        draw=obeBlue, line width=0.8pt, minimum height=0.85cm,
        align=center, inner sep=2pt
    },
    ket/.style={font=\scriptsize, text=obeGray, align=center}
]
    \node[bit, fill=obeAmberLight, minimum width=1.6cm] (b76)
        at (0,0) {\textbf{7}\quad\textbf{6}};
    \node[bit, fill=obeBlueLight, minimum width=2.4cm,
        right=0.12cm of b76] (b53)
        {\textbf{5}\quad\textbf{4}\quad\textbf{3}};
    \node[bit, fill=obeTealLight, minimum width=2.4cm,
        right=0.12cm of b53] (b20)
        {\textbf{2}\quad\textbf{1}\quad\textbf{0}};

    \node[below=0.06cm of b76, font=\bfseries] {\code{mod}};
    \node[below=0.06cm of b53, font=\bfseries] {\code{reg/opcode}};
    \node[below=0.06cm of b20, font=\bfseries] {\code{r/m}};

    \node[ket, below=0.55cm of b76] {2 bit\\cara pengambilan};
    \node[ket, below=0.55cm of b53] {3 bit\\register kedua};
    \node[ket, below=0.55cm of b20] {3 bit\\register / alamat};
\end{tikzpicture}
\end{center}

\vspace{0.5em}
{\small Contoh byte ModR/M \hex{D8} pada \code{MOV AX, BX} ditulis
\bin{11 011 000}: \code{mod} $= \bin{11}$, \code{reg}
$= \bin{011}$ (\reg{BX}), \code{r/m} $= \bin{000}$ (\reg{AX}).
\code{mod} $= \bin{11}$ berarti kedua operand berada di register.\par}
\Sumber{Bab VII, Struktur Byte ModR/M}
\end{frame}

% ---------------------------------------------------------
\begin{frame}{Arti Nilai Field \code{mod}}
\footnotesize Dua bit \code{mod} menentukan letak operand dan panjang
\eng{displacement} yang menyertainya.

\vspace{0.35em}
\begin{center}
\begin{tabular}{@{}ccp{2.2cm}p{7.4cm}@{}}
    \toprule
    \textbf{\code{mod}} & \textbf{Biner} & \textbf{Displacement} &
    \textbf{Arti} \\
    \midrule
    \dec{0} & \bin{00} & tidak ada &
    Operand memori; alamat langsung dari \code{r/m} \\
    \addlinespace
    \dec{1} & \bin{01} & 1 byte (disp8) &
    Operand memori; alamat ditambah \eng{displacement} 8 bit \\
    \addlinespace
    \dec{2} & \bin{10} & 2 byte (disp16) &
    Operand memori; alamat ditambah \eng{displacement} 16 bit \\
    \addlinespace
    \dec{3} & \bin{11} & tidak ada &
    Operand register; \code{r/m} menyatakan nomor register \\
    \bottomrule
\end{tabular}
\end{center}

\vspace{0.45em}
{\color{obeGray}Memakai \code{mod} $= \bin{01}$ lebih hemat satu byte
daripada $\bin{10}$ untuk alamat yang sama, asalkan pergeserannya muat
dalam 8 bit.\par}
\end{frame}

% ---------------------------------------------------------
\begin{frame}{Arti Nilai Field \code{r/m}}
\footnotesize Field \code{r/m} berperan ganda: penentu alamat efektif
bila \code{mod} $\neq \bin{11}$, dan nomor register bila
\code{mod} $= \bin{11}$.

\vspace{0.25em}
\begin{columns}[t]
\begin{column}{0.575\textwidth}
{\color{obeBlue}\bfseries Mode memori (\code{mod} $\neq \bin{11}$)}\\
\vspace{0.15em}
\begin{tabular}{@{}ccp{2.3cm}p{2.4cm}@{}}
    \toprule
    \textbf{r/m} & \textbf{Biner} & \code{mod}$=$\bin{00} &
    \code{mod}$=$\bin{01}/\bin{10} \\
    \midrule
    \dec{0} & \bin{000} & \asm{[BX+SI]} & \asm{[BX+SI]}+disp \\
    \dec{1} & \bin{001} & \asm{[BX+DI]} & \asm{[BX+DI]}+disp \\
    \dec{2} & \bin{010} & \asm{[BP+SI]} & \asm{[BP+SI]}+disp \\
    \dec{3} & \bin{011} & \asm{[BP+DI]} & \asm{[BP+DI]}+disp \\
    \dec{4} & \bin{100} & \asm{[SI]} & \asm{[SI]}+disp \\
    \dec{5} & \bin{101} & \textbf{langsung} & \asm{[DI]}+disp \\
    \dec{6} & \bin{110} & \asm{[BP]} & \asm{[BP]}+disp \\
    \dec{7} & \bin{111} & \asm{[BX]} & \asm{[BX]}+disp \\
    \bottomrule
\end{tabular}
\end{column}
\begin{column}{0.40\textwidth}
{\color{obeTeal}\bfseries Mode register (\code{mod} $= \bin{11}$)}\\
\vspace{0.15em}
\begin{tabular}{@{}ccc@{}}
    \toprule
    \textbf{r/m} & \textbf{Reg 16} & \textbf{Reg 8} \\
    \midrule
    \dec{0} & \reg{AX} & \reg{AL} \\
    \dec{1} & \reg{CX} & \reg{CL} \\
    \dec{2} & \reg{DX} & \reg{DL} \\
    \dec{3} & \reg{BX} & \reg{BL} \\
    \dec{4} & \reg{SP} & \reg{AH} \\
    \dec{5} & \reg{BP} & \reg{CH} \\
    \dec{6} & \reg{SI} & \reg{DH} \\
    \dec{7} & \reg{DI} & \reg{BH} \\
    \bottomrule
\end{tabular}
\end{column}
\end{columns}
\Sumber{Bab VII, Arti Nilai r/m}
\end{frame}

% ---------------------------------------------------------
\begin{frame}{Dua Pengecualian dan Urutan Nomor Register}
\begin{alertblock}{Dua pengecualian yang sering menjebak}
\small
\begin{enumerate}
    \item Bila \code{mod} $= \bin{00}$ dan \code{r/m} $= \bin{101}$,
          alamat efektifnya bukan \asm{[DI]} melainkan \textbf{alamat
          langsung 16 bit}.
    \item Bila \code{mod} $\neq \bin{11}$ dan \code{r/m} $= \bin{110}$,
          alamat efektifnya adalah \asm{[BP]}.
\end{enumerate}
Kedua aturan ini di luar pola umum dan harus dihafal.
\end{alertblock}

\vspace{0.45em}
\begin{block}{Urutan nomor register 8086}
\small \reg{AX}, \reg{CX}, \reg{DX}, \reg{BX}, \reg{SP}, \reg{BP},
\reg{SI}, \reg{DI} --- mengikuti \eng{encoding} asli Intel dan
berlaku juga pada \asm{PUSH} serta \asm{POP}.
\end{block}

\vspace{0.45em}
{\footnotesize\color{obeGray}\code{mod} $= \bin{01}$ dan $\bin{10}$
memakai tabel alamat yang sama; hanya panjang
\eng{displacement}-nya berbeda. Contoh: \asm{ADD AX, [SI+4]} memakai
\code{mod} $= \bin{01}$, \code{r/m} $= \bin{100}$, ModR/M
\hex{44}.\par}
\end{frame}

% ---------------------------------------------------------
\begin{frame}{Peran Displacement dan Immediate}
\small Kedua field terakhir sama-sama berupa bilangan, tetapi perannya
berbeda.
\vspace{0.3em}
\begin{description}
    \item[Displacement] Penambah pada alamat. Muncul hanya bila operand
          berupa memori dan \code{mod} bernilai \bin{01} atau
          \bin{10}. Pada \asm{ADD AX, [SI+4]}, nilai \hex{4} adalah
          \eng{displacement} sepanjang 8 atau 16 bit sesuai
          \code{mod}.

    \vspace{0.35em}
    \item[Immediate] Data yang dioperasikan. Muncul bila operand berupa
          konstanta, misalnya \asm{MOV AX, 1234h} atau
          \asm{ADD BX, 10}. Panjangnya 8 atau 16 bit sesuai ukuran
          operand, dan nilainya selalu \eng{little-endian}.
\end{description}

\vspace{0.4em}
\begin{exampleblock}{Keduanya dapat hadir bersamaan}
\small \asm{MOV WORD PTR [BX+SI+1234h], 5678h} memuat
\eng{displacement} dan \eng{immediate} sekaligus, sehingga panjang
totalnya enam byte. Urutan field tetap: \eng{displacement} selalu
mendahului \eng{immediate}.
\end{exampleblock}
\end{frame}

% ---------------------------------------------------------
\begin{frame}{Pengaruh Format terhadap Ukuran dan Kecepatan}
\small
\begin{center}
\begin{tabular}{@{}p{4.2cm}cp{6.8cm}@{}}
    \toprule
    \textbf{Bentuk operand} & \textbf{Panjang} & \textbf{Keterangan} \\
    \midrule
    Register ke register & 2 byte &
    Paling ringkas; tanpa \eng{displacement} dan tanpa
    \eng{immediate} \\
    \addlinespace
    Immediate 8 bit & 2--3 byte & Nilai kecil; satu byte
    \eng{immediate} \\
    \addlinespace
    Immediate 16 bit & 3--4 byte & Nilai besar; dua byte
    \eng{immediate} \\
    \addlinespace
    Memori \asm{[SI]} & 2--3 byte & \code{mod} $= \bin{00}$; tanpa
    \eng{displacement} \\
    \addlinespace
    Memori \asm{[SI+4]} & 3 byte & \code{mod} $= \bin{01}$;
    \eng{displacement} 8 bit \\
    \addlinespace
    Memori \asm{[BX+SI+1234h]} & 4 byte & \code{mod} $= \bin{10}$;
    \eng{displacement} 16 bit \\
    \bottomrule
\end{tabular}
\end{center}

\vspace{0.3em}
{\footnotesize\color{obeGray}Semakin banyak field yang hadir, semakin
panjang instruksi dan semakin banyak byte yang harus diambil melalui
\eng{prefetch queue}. Instruksi tiga byte umumnya selesai lebih cepat
daripada instruksi enam byte.\par}
\Sumber{Bab VII, Pengaruh Format terhadap Ukuran dan Kecepatan}
\end{frame}

% ---------------------------------------------------------
\begin{frame}{Rangkuman Bagian 2}
\small
\begin{itemize}
    \item Byte ModR/M terbagi menjadi \code{mod} (2 bit),
          \code{reg/opcode} (3 bit), dan \code{r/m} (3 bit).

    \vspace{0.25em}
    \item Nilai \code{mod} menentukan letak operand dan panjang
          \eng{displacement}: \bin{00} tanpa, \bin{01} disp8,
          \bin{10} disp16, dan \bin{11} operand register.

    \vspace{0.25em}
    \item \code{r/m} berperan ganda: penentu alamat efektif pada mode
          memori dan nomor register pada \code{mod} $= \bin{11}$.

    \vspace{0.25em}
    \item Dua pengecualian: \code{mod} $= \bin{00}$ dengan
          \code{r/m} $= \bin{101}$ berarti alamat langsung, sedangkan
          \code{r/m} $= \bin{110}$ berarti \asm{[BP]}.

    \vspace{0.25em}
    \item \eng{Displacement} menggeser alamat, \eng{immediate} menjadi
          operand. Bentuk operand menentukan ukuran kode dan jumlah
          byte yang diambil prosesor.
\end{itemize}
\end{frame}

% =========================================================
% BAGIAN 3
% =========================================================
\PemisahBagian{Encoding Manual dan Disassembler}

% ---------------------------------------------------------
\begin{frame}{Lima Langkah Encoding Manual}
\small Mengubah satu instruksi assembly menjadi byte kode mesin:
\vspace{0.3em}
\begin{enumerate}
    \item Tentukan operasi dan bentuk operandnya: register, memori,
          atau nilai segera --- serta ukurannya (8 atau 16 bit).

    \vspace{0.2em}
    \item Cari \eng{opcode} dasar yang sesuai. Perhatikan arah
          perpindahan data, karena \hex{8B} dan \hex{89} memiliki arti
          berlawanan.

    \vspace{0.2em}
    \item Hitung byte ModR/M bila instruksi melibatkan register atau
          memori: susun field \code{mod}, \code{reg}, dan \code{r/m},
          lalu gabungkan menjadi satu byte.

    \vspace{0.2em}
    \item Tambahkan \eng{displacement} bila operand memori bergeser,
          dan \eng{immediate} bila operandnya berupa angka.

    \vspace{0.2em}
    \item Tuliskan seluruh byte berurutan menurut alamat, dengan setiap
          nilai 16 bit dalam urutan \eng{little-endian}.
\end{enumerate}

\vspace{0.25em}
{\footnotesize\color{obeGray}Pekerjaan manual inilah yang dijalankan
assembler untuk setiap baris kode sumber.\par}
\Sumber{Bab VII, Langkah-langkah Encoding Manual}
\end{frame}

% ---------------------------------------------------------
\begin{frame}{Contoh Encoding Byte demi Byte}
\footnotesize
\begin{center}
\begin{tabular}{@{}p{3.1cm}p{1.7cm}p{7.6cm}@{}}
    \toprule
    \textbf{Instruksi} & \textbf{Byte} & \textbf{Rincian} \\
    \midrule
    \asm{MOV AX, BX} & \hex{89 D8} &
    Opcode \hex{89} = ``r/m16 $\leftarrow$ r16''. ModR/M \hex{D8}
    $= \bin{11 011 000}$: \code{mod} $= \bin{11}$, \code{reg}
    $= \bin{011}$ (\reg{BX}), \code{r/m} $= \bin{000}$ (\reg{AX}) \\
    \addlinespace
    \asm{MOV AL, 5} & \hex{B0 05} &
    Opcode \hex{B0} = ``r8 $\leftarrow$ imm8''. Nomor \reg{AL} =
    \dec{0}, jadi \hex{B0} $+ \dec{0} = \hex{B0}$. Byte kedua
    \hex{05} adalah \eng{immediate} 8 bit \\
    \addlinespace
    \asm{MOV AX, 1234h} & \hex{B8 34 12} &
    Opcode \hex{B8} = ``r16 $\leftarrow$ imm16''. Nomor \reg{AX} =
    \dec{0}. \eng{Immediate} \hex{1234} disimpan
    \eng{little-endian}: \hex{34} lalu \hex{12} \\
    \addlinespace
    \asm{ADD AX, [SI+4]} & \hex{03 44 04} &
    Opcode \hex{03}. ModR/M \hex{44} $= \bin{01 000 100}$:
    \code{mod} $= \bin{01}$ (disp8), \code{reg} $= \bin{000}$
    (\reg{AX}), \code{r/m} $= \bin{100}$ (\asm{[SI]}). Byte ketiga
    \hex{04} adalah \eng{displacement} \\
    \addlinespace
    \asm{INT 21h} & \hex{CD 21} &
    Opcode \hex{CD} = ``panggil interupsi''. Tidak memakai ModR/M;
    byte \hex{21} adalah \eng{immediate} 8 bit \\
    \bottomrule
\end{tabular}
\end{center}
\end{frame}

% ---------------------------------------------------------
\begin{frame}{Tabel Ringkas Opcode: MOV dan Aritmatika}
\footnotesize
\begin{center}
\begin{tabular}{@{}p{1.8cm}p{2.2cm}p{6.4cm}c@{}}
    \toprule
    \textbf{Instruksi} & \textbf{Opcode} & \textbf{Bentuk operand} &
    \textbf{Byte} \\
    \midrule
    \asm{MOV} & \hex{88}--\hex{8B} & Antar \code{r/m} dan register &
    2--4 \\
    \addlinespace
    \asm{MOV} & \hex{B0}--\hex{B7} & Register 8 bit,
    \eng{immediate} 8 bit & 2 \\
    \addlinespace
    \asm{MOV} & \hex{B8}--\hex{BF} & Register 16 bit,
    \eng{immediate} 16 bit & 3 \\
    \addlinespace
    \asm{MOV} & \hex{C6}/\hex{C7} & \code{r/m}, \eng{immediate} &
    3--6 \\
    \addlinespace
    \asm{ADD} & \hex{00}--\hex{03} & Antar \code{r/m} dan register &
    2--4 \\
    \addlinespace
    \asm{ADD} & \hex{04}/\hex{05} & \reg{AL}/\reg{AX} dengan
    \eng{immediate} & 2--3 \\
    \addlinespace
    \asm{SUB} & \hex{28}--\hex{2D} & Pola sama dengan \asm{ADD} &
    2--4 \\
    \addlinespace
    \asm{CMP} & \hex{38}--\hex{3D} & Pembandingan; hanya mengubah
    \eng{flag} & 2--4 \\
    \bottomrule
\end{tabular}
\end{center}

\vspace{0.25em}
{\footnotesize\color{obeGray}\hex{00}--\hex{03} berarti empat
\eng{opcode} berurutan: \hex{00} untuk 8 bit arah \code{r/m}
$\leftarrow$ \code{reg}, \hex{01} untuk 16 bit, lalu \hex{02} dan
\hex{03} untuk arah sebaliknya.\par}
\end{frame}

% ---------------------------------------------------------
\begin{frame}{Tabel Ringkas Opcode: Lompatan, Interupsi, dan Stack}
\footnotesize
\begin{center}
\begin{tabular}{@{}p{1.8cm}p{2.2cm}p{6.4cm}c@{}}
    \toprule
    \textbf{Instruksi} & \textbf{Opcode} & \textbf{Bentuk operand} &
    \textbf{Byte} \\
    \midrule
    \asm{JMP} & \hex{EB} & Lompat pendek, \eng{rel8} & 2 \\
    \addlinespace
    \asm{JMP} & \hex{E9} & Lompat dekat, \eng{rel16} & 3 \\
    \addlinespace
    \asm{CALL} & \hex{E8} & Panggil dekat, \eng{rel16} & 3 \\
    \addlinespace
    \asm{INT} & \hex{CC} & \code{INT 3}; titik henti
    \eng{debugger} & 1 \\
    \addlinespace
    \asm{INT} & \hex{CD} & Interupsi dengan nomor 8 bit & 2 \\
    \addlinespace
    \asm{PUSH} & \hex{50}--\hex{57} & Dorong register 16 bit ke
    \eng{stack} & 1 \\
    \addlinespace
    \asm{POP} & \hex{58}--\hex{5F} & Ambil register 16 bit dari
    \eng{stack} & 1 \\
    \bottomrule
\end{tabular}
\end{center}

\vspace{0.4em}
\begin{block}{Pola kelompok}
\small Menghafal seluruh \eng{opcode} 8086 tidak diperlukan; mengenali
pola kelompoknya jauh lebih berguna.
\end{block}
\Sumber{Bab VII, Tabel Ringkas Opcode Instruksi Umum}
\end{frame}

% ---------------------------------------------------------
\begin{frame}[fragile]{Membalik Kode Mesin dengan Disassembler}
\small \term{Disassembler} melakukan pekerjaan kebalikan assembler:
membaca byte kode mesin dan menampilkannya kembali sebagai
\eng{mnemonic}. \file{DEBUG} menyediakan perintah \code{U}
(\eng{unassemble}) untuk keperluan ini.

\vspace{0.25em}
\begin{keluaran}
-u 100 10C
13F2:0100 B80500   MOV   AX,0005
13F2:0103 BB0100   MOV   BX,0001
13F2:0106 8908     MOV   [BX+SI],CX
13F2:0108 03C3     ADD   AX,BX
13F2:010A CD21     INT   21
13F2:010C 90       NOP
\end{keluaran}

\vspace{0.2em}
{\footnotesize\color{obeGray}Setiap baris memuat tiga kolom: alamat,
byte kode mesin, lalu \eng{mnemonic}. Baris pertama menempati tiga
byte, sedangkan \asm{ADD AX, BX} hanya dua byte.\par}

\vspace{0.35em}
\begin{alertblock}{Keterbatasan yang harus diingat}
\small Disassembler tidak dapat membedakan kode dari data. Hasilnya
harus selalu diperiksa terhadap titik masuk dan alur kendali program.
\end{alertblock}
\end{frame}

% ---------------------------------------------------------
\begin{frame}{Hubungan dengan Translasi Program}
\small
\begin{block}{Encoding manual adalah inti proses translasi}
\begin{itemize}
    \item Assembler menjalankan lima langkah \eng{encoding} yang sama
          untuk setiap baris kode sumber, tetapi jauh lebih cepat dan
          teliti.
    \item Direktif, tabel simbol, dan penanganan label dibahas pada
          Bab IX.
    \item Dengan memahami format instruksi lebih dahulu, setiap
          keputusan assembler dapat ditelusuri sampai byte demi byte.
\end{itemize}
\end{block}

\vspace{0.5em}
\begin{center}
\begin{tikzpicture}[
    kotak/.style={
        draw=obeBlue, line width=0.8pt, fill=obeBlueLight,
        rounded corners=3pt, minimum width=3.0cm, minimum height=0.95cm,
        align=center, font=\small, inner sep=3pt
    },
    panah/.style={-{Stealth[length=2mm]}, line width=0.9pt, obeTeal}
]
    \node[kotak] (s) {\file{.asm}\\{\scriptsize kode sumber}};
    \node[kotak, right=0.9cm of s, fill=obeTealLight, draw=obeTeal] (a)
        {Assembler\\{\scriptsize lima langkah}};
    \node[kotak, right=0.9cm of a, fill=obeAmberLight, draw=obeAmber] (b)
        {\hex{89 D8}\\{\scriptsize byte kode mesin}};
    \draw[panah] (s) -- (a);
    \draw[panah] (a) -- (b);
\end{tikzpicture}
\end{center}
\end{frame}

% ---------------------------------------------------------
\begin{frame}{Rangkuman Bagian 3}
\small
\begin{itemize}
    \item \eng{Encoding} manual ditempuh lewat lima langkah: tentukan
          operasi dan operand, cari \eng{opcode}, susun ModR/M,
          tambahkan \eng{displacement}/\eng{immediate}, lalu tulis
          \eng{little-endian}.

    \vspace{0.25em}
    \item Sebagian \eng{opcode} memakai tiga bit \code{reg} sebagai
          penambah, seperti \hex{B0} dan \hex{B8} yang nomor
          register-nya disisipkan pada bit terendah.

    \vspace{0.25em}
    \item Instruksi yang tidak melibatkan register atau memori sama
          sekali tidak memerlukan byte ModR/M, misalnya
          \asm{INT 21h}.

    \vspace{0.25em}
    \item Perintah \code{U} pada \file{DEBUG} mencetak alamat, byte
          kode mesin, dan \eng{mnemonic} hasil terjemahan.

    \vspace{0.25em}
    \item Disassembler tidak membedakan kode dari data; hasilnya harus
          diperiksa terhadap titik masuk dan alur kendali program.
\end{itemize}
\end{frame}

% =========================================================
% CONTOH PROGRAM
% =========================================================
\begin{frame}[fragile]{Contoh 1: Kode Sumber \file{opcode-demo.asm}}
\small Lima instruksi dengan bentuk operand berbeda, siap dibongkar
dengan perintah \code{U}.

\vspace{0.2em}
\begin{kodeasm}
.MODEL SMALL
.DATA
    pesan DB 'Format instruksi mesin', 0Dh, 0Ah, '$'
.CODE
START:
    MOV AX, @DATA
    MOV DS, AX
    MOV AX, 0005h      ; B8 05 00
    MOV BX, 0001h      ; BB 01 00
    ADD AX, BX         ; 03 C3
    MOV CX, AX         ; 89 C1
    SUB AX, 0002h      ; 2D 02 00
    ...                ; cetak pesan dan keluar dari program
END START
\end{kodeasm}

\vspace{0.15em}
{\scriptsize\color{obeGray}Bagian pencetakan pesan dan keluar program
(\asm{MOV AH, 09h}, \asm{LEA DX, pesan}, \asm{INT 21h},
\asm{MOV AH, 4Ch}) dihilangkan agar ringkas; berkas lengkapnya ada pada
\file{code/asm/opcode-demo.asm}.\par}
\end{frame}

% ---------------------------------------------------------
\begin{frame}[fragile]{Contoh 1: Hasil Bongkar dan Perbandingannya}
\footnotesize
\begin{keluaran}
-u
0B2A:0005 B80500   MOV   AX,0005
0B2A:0008 BB0100   MOV   BX,0001
0B2A:000B 03C3     ADD   AX,BX
0B2A:000D 89C1     MOV   CX,AX
0B2A:000F 2D0200   SUB   AX,0002
\end{keluaran}

\vspace{0.25em}
\begin{center}
\begin{tabular}{@{}p{2.9cm}p{2.1cm}cp{6.2cm}@{}}
    \toprule
    \textbf{Kode sumber} & \textbf{Kode mesin} & \textbf{Byte} &
    \textbf{Field yang tampak} \\
    \midrule
    \asm{MOV AX, 0005h} & \hex{B8 05 00} & 3 &
    Opcode penambah register, imm.\ 16 bit \\
    \addlinespace
    \asm{MOV BX, 0001h} & \hex{BB 01 00} & 3 &
    Nomor \reg{BX} $= \dec{3}$, jadi \hex{B8} $+ \dec{3}$ \\
    \addlinespace
    \asm{ADD AX, BX} & \hex{03 C3} & 2 &
    ModR/M \bin{11 000 011}: \code{reg} $= \reg{AX}$,
    \code{r/m} $= \reg{BX}$ \\
    \addlinespace
    \asm{MOV CX, AX} & \hex{89 C1} & 2 &
    Opcode \hex{89} arah \code{r/m} $\leftarrow$ \code{reg} \\
    \addlinespace
    \asm{SUB AX, 0002h} & \hex{2D 02 00} & 3 &
    Opcode AX dengan immediate 16 bit \\
    \bottomrule
\end{tabular}
\end{center}
\end{frame}

% ---------------------------------------------------------
\begin{frame}[fragile]{Contoh 2: Kode Sumber \file{ukuran-instruksi.asm}}
\small Dua penulisan yang menghasilkan keadaan register sama, tetapi
menempati jumlah byte berbeda.

\vspace{0.2em}
\begin{kodeasm}
.MODEL TINY
.CODE
ORG 100h
START:
    MOV AX, 0001h      ; bentuk A, 3 byte
    MOV BX, 0001h      ; bentuk A, 3 byte
    MOV BX, 0001h      ; bentuk B, 3 byte
    MOV AX, BX         ; bentuk B, 2 byte
    ...
pesan DB 'Dua bentuk, hasil sama, ukuran berbeda.', 0Dh, 0Ah, '$'
END START
\end{kodeasm}

\vspace{0.2em}
{\footnotesize\color{obeGray}Program berbentuk \file{.COM}, sehingga
seluruh instruksi dimulai pada offset \hex{100} dan \reg{DS} tidak
perlu diinisialisasi.\par}
\end{frame}

% ---------------------------------------------------------
\begin{frame}[fragile]{Contoh 2: Hasil Bongkar dan Jumlah Byte}
\footnotesize
\begin{keluaran}
-u 100 10A
13F2:0100 B80100   MOV   AX,0001   ; bentuk A, 3 byte
13F2:0103 BB0100   MOV   BX,0001   ; bentuk A, 3 byte
13F2:0106 BB0100   MOV   BX,0001   ; bentuk B, 3 byte
13F2:0109 89D8     MOV   AX,BX     ; bentuk B, 2 byte
\end{keluaran}

\vspace{0.25em}
\begin{center}
\begin{tabular}{@{}cp{4.6cm}p{3.0cm}c@{}}
    \toprule
    \textbf{Bentuk} & \textbf{Instruksi} & \textbf{Kode mesin} &
    \textbf{Byte} \\
    \midrule
    A & \asm{MOV AX, 0001h} & \hex{B8 01 00} & 3 \\
    \addlinespace
    A & \asm{MOV BX, 0001h} & \hex{BB 01 00} & 3 \\
    \addlinespace
    \multicolumn{3}{@{}l}{\textbf{Jumlah bentuk A}} & \textbf{6} \\
    \midrule
    B & \asm{MOV BX, 0001h} & \hex{BB 01 00} & 3 \\
    \addlinespace
    B & \asm{MOV AX, BX} & \hex{89 D8} & 2 \\
    \addlinespace
    \multicolumn{3}{@{}l}{\textbf{Jumlah bentuk B}} & \textbf{5} \\
    \bottomrule
\end{tabular}
\end{center}

\vspace{0.2em}
{\scriptsize\color{obeGray}Selisih satu byte berasal dari
\eng{immediate} 16 bit kedua yang tidak lagi diperlukan. Bentuk B belum
tentu lebih baik: ia bergantung pada nilai register lain dan lebih
sulit dibaca.\par}
\end{frame}

% =========================================================
% PRAKTIKUM
% =========================================================
\begin{frame}{Praktikum: Analisis Opcode Manual dan Disassembler}
\begin{block}{Judul praktikum}
\scriptsize Analisis \eng{opcode} manual dan \eng{disassembler}
\quad\textbullet\quad \textbf{Sub-CPMK:} 2.1 \quad\textbullet\quad
\textbf{Waktu:} 100 menit \quad\textbullet\quad \textbf{Bentuk:}
individu \quad\textbullet\quad \textbf{Tugas 2}
\end{block}

\vspace{0.25em}
\begin{columns}[t]
\begin{column}{0.48\textwidth}
{\color{obeBlue}\bfseries Alat dan bahan}
\begin{itemize}\footnotesize
    \item Komputer bersistem operasi Windows
    \item Emu8086, GUI Turbo Assembler (TASM/TLINK), atau DOSBox
          dengan \file{DEBUG}
    \item Berkas \file{opcode-demo.asm},
          \file{ukuran-instruksi.asm}, dan lembar kerja
          \eng{encoding}
\end{itemize}
\end{column}
\begin{column}{0.48\textwidth}
{\color{obeBlue}\bfseries Langkah kerja}
\begin{enumerate}\footnotesize
    \item Isi lembar kerja \eng{encoding} lima instruksi secara
          manual
    \item Rakit \file{opcode-demo.asm}; jalankan \code{U} pada
          \file{DEBUG}
    \item Cocokkan tiap baris hasil bongkar dengan kode sumber
    \item Rakit \file{ukuran-instruksi.asm}; hitung byte bentuk A
          dan B
    \item Ubah \asm{MOV AX, 1234h} menjadi \asm{MOV AX, 3412h};
          jelaskan \eng{little-endian}-nya
\end{enumerate}
\end{column}
\end{columns}

\vspace{0.15em}
{\scriptsize\textbf{Dikumpulkan:} lembar kerja \eng{encoding}
terisi, keluaran perintah \code{U}, perhitungan byte bentuk A dan B,
jawaban penelusuran \eng{immediate}, dan jurnal praktikum.\par}
\end{frame}

% =========================================================
% LATIHAN
% =========================================================
\begin{frame}{Latihan}
\small
\begin{enumerate}
    \item Jelaskan hubungan antara \eng{mnemonic} dan \eng{opcode}.
          Mengapa programmer cukup menulis \asm{MOV AX, BX} tanpa
          menentukan sendiri byte \hex{89 D8}?

    \vspace{0.25em}
    \item Instruksi berbyte \hex{C7 80 34 12 78 56}. Tentukan
          \eng{opcode} dan operasinya, nilai biner \code{mod},
          \code{reg}, dan \code{r/m}, serta panjang dan isi
          \eng{displacement} dan \eng{immediate}.

    \vspace{0.25em}
    \item Tuliskan byte kode mesin untuk \asm{MOV DX, 00FFh},
          \asm{ADD AL, 3}, dan \asm{MOV [BX], AL} beserta langkah
          perhitungannya.

    \vspace{0.25em}
    \item Jelaskan perbedaan peran \eng{displacement} dan
          \eng{immediate}. Kapan keduanya muncul, dan dapatkah
          keduanya muncul sekaligus?

    \vspace{0.25em}
    \item Dua program berakhir sama. Program pertama memakai
          \asm{MOV AX, 0001h} dan \asm{MOV BX, 0001h}; program kedua
          memakai \asm{MOV BX, 0001h} dan \asm{MOV AX, BX}. Hitung
          jumlah byte keduanya dan sebutkan field penyebab
          selisihnya.
\end{enumerate}
\end{frame}

% =========================================================
% RANGKUMAN
% =========================================================
\begin{frame}{Rangkuman}
\small
\begin{itemize}
    \item \eng{Mnemonic} assembly hanyalah simbol bagi kode mesin.
          Satu mnemonic memiliki banyak bentuk \eng{encoding} karena
          jenis operand, ukuran operand, arah perpindahan data, dan
          mode pengalamatan.

    \vspace{0.2em}
    \item Format instruksi 8086 tersusun dari lima field berurutan
          tetap: awalan, \eng{opcode}, ModR/M, \eng{displacement},
          dan \eng{immediate}. Panjang totalnya 1 sampai 6 byte.

    \vspace{0.2em}
    \item Byte ModR/M terbagi menjadi \code{mod} (2 bit),
          \code{reg/opcode} (3 bit), dan \code{r/m} (3 bit). Nilai
          \code{mod} menentukan letak operand dan panjang
          \eng{displacement}.

    \vspace{0.2em}
    \item \eng{Displacement} menggeser alamat efektif dan hanya muncul
          pada operand memori; \eng{immediate} adalah konstanta yang
          menjadi operand. Keduanya dapat muncul bersamaan.

    \vspace{0.2em}
    \item Setiap nilai 16 bit disimpan \eng{little-endian}, dan bentuk
          operand mempengaruhi ukuran kode serta jumlah byte yang
          diambil prosesor.

    \vspace{0.2em}
    \item Disassembler, misalnya perintah \code{U} pada \file{DEBUG},
          membalik kode mesin menjadi \eng{mnemonic}; hasilnya harus
          diperiksa karena disassembler tidak membedakan kode dari
          data.
\end{itemize}
\end{frame}

% ---------------------------------------------------------
\begin{frame}{Cek Pemahaman}
\small
\begin{enumerate}
    \item Urutan field instruksi 8086 dari kiri ke kanan adalah \dots
          \begin{enumerate}\footnotesize
              \item \eng{opcode}, awalan, ModR/M, \eng{immediate},
                    \eng{displacement}
              \item awalan, \eng{opcode}, ModR/M, \eng{displacement},
                    \eng{immediate}
              \item ModR/M, \eng{opcode}, awalan, \eng{displacement},
                    \eng{immediate}
              \item awalan, ModR/M, \eng{opcode}, \eng{immediate},
                    \eng{displacement}
          \end{enumerate}

    \vspace{0.2em}
    \item Byte ModR/M \hex{D8} memiliki nilai \code{mod} sebesar \dots
          \begin{enumerate}\footnotesize
              \item \bin{00}, operand memori tanpa \eng{displacement}
              \item \bin{01}, operand memori dengan disp8
              \item \bin{10}, operand memori dengan disp16
              \item \bin{11}, operand berupa register
          \end{enumerate}

    \vspace{0.2em}
    \item \asm{MOV AX, 1234h} dikodekan \hex{B8 34 12} dan bukan
          \hex{B8 12 34} karena \dots
          \begin{enumerate}\footnotesize
              \item 8086 menyimpan nilai 16 bit secara
                    \eng{little-endian}
              \item \eng{opcode} \hex{B8} hanya melayani nilai kecil
              \item urutan byte ditentukan oleh mode pengalamatan
              \item byte rendah disimpan pada \reg{AH}
          \end{enumerate}
\end{enumerate}
\end{frame}

% =========================================================
% PENUTUP
% =========================================================
\begin{frame}{Penutup}
\small
\begin{block}{Sebelum pertemuan berikutnya}
\begin{itemize}
    \item Selesaikan Tugas 2: lembar kerja \eng{encoding} terisi,
          keluaran perintah \code{U}, dan perhitungan byte dua bentuk
          penulisan.
    \item Latih kembali pemecahan byte ModR/M dan urutan
          \eng{little-endian} sebelum mempelajari mode pengalamatan.
    \item Baca \textbf{Bab VIII} --- Mode Pengalamatan.
\end{itemize}
\end{block}

\vspace{0.4em}
\begin{center}
{\color{obeBlue}\sffamily\large\bfseries Terima kasih\par}
\vspace{0.2em}
{\color{obeGray}\sffamily\small \NamaMK\ (\KodeMK) \textbullet\
\SemesterTeks\par}
\end{center}
\end{frame}

\end{document}
